\documentclass[10pt,aspectratio=169]{beamer}
\usepackage{poly}
\usepackage[useregional]{datetime2}
\usepackage{graphicx}  % for images
\usepackage[table]{xcolor}    % for color with table support
\usepackage{minted}    % for better code highlighting
\usepackage[T1]{fontenc}  % for better font encoding
\usepackage{inconsolata}  % fallback monospace font

% Minted configuration
\setminted{
    frame=none,
    fontsize=\small,
    linenos=true,
    numbersep=12pt,
    bgcolor=white,
    style=friendly,
    autogobble,
    breaklines,
    baselinestretch=1.2
}

% Add slide numbers to the footer
\setbeamertemplate{footline}[frame number]
\setbeamertemplate{navigation symbols}{}  % Remove navigation symbols

% Define a centered minted environment
\newenvironment{centeredcode}{\begin{center}}{\end{center}}

\title{Operating System (IF2223)}
\subtitle{Memory Management 1}
\author{Martin C.T. Manullang, Eko Dwi Nugroho, Ilham Firman Ashari}
\institute[COMP]{Teknik Informatika, ITERA}
\date{2025}

\begin{document}
\maketitle

\section{Pendahuluan}

% --- Frame: Memory in Operating Systems ---
\begin{frame}{Memori dalam Sistem Operasi}
    \framesubtitle{Komponen Utama dalam Operasi Sistem}
    \begin{itemize}
        \item Memori adalah komponen \textbf{utama} dalam operasi sistem komputer modern
        \item Terdiri dari array byte yang besar, masing-masing dengan alamat tersendiri
        \item CPU berinteraksi dengan memori untuk:
        \begin{itemize}
            \item Mengambil instruksi
            \item Menyimpan dan mengambil operand
            \item Menulis hasil operasi
        \end{itemize}
        \item Unit memori hanya melihat aliran alamat
        \item Memori tidak mengetahui bagaimana alamat dihasilkan atau tujuannya
    \end{itemize}
\end{frame}

% --- Frame: Instruction Execution Cycle ---
\begin{frame}{Siklus Eksekusi Instruksi}
    \framesubtitle{Peran Memori}
    \begin{columns}[T]
        \column{0.48\textwidth}
            \textbf{Siklus Instruksi Tipikal:}
            \begin{enumerate}
                \item Ambil instruksi dari memori
                \item Decode instruksi
                \item Ambil operand dari memori (jika diperlukan)
                \item Eksekusi instruksi
                \item Simpan hasil ke memori (jika diperlukan)
            \end{enumerate}
        
        \column{0.48\textwidth}
            \begin{alertblock}{Konsep Kunci}
                Untuk eksekusi yang tepat, instruksi dan data harus dapat diakses dalam memori saat dibutuhkan.
            \end{alertblock}
            
            \begin{block}{Fokus Manajemen Memori}
                Fokus kita adalah pada urutan alamat memori yang dihasilkan oleh program yang berjalan, bukan bagaimana alamat tersebut dihasilkan.
            \end{block}
    \end{columns}
\end{frame}

% --- Frame: Siklus Eksekusi Instruksi - Dasar ---
\begin{frame}{Siklus Eksekusi Instruksi}
    \framesubtitle{Konsep Dasar}
    
    \begin{block}{Definisi}
        Serangkaian langkah yang diambil CPU untuk mengeksekusi setiap instruksi program.
    \end{block}
    
    \begin{columns}[T]
        \column{0.48\textwidth}
            \textbf{Langkah-langkah Utama:}
            \begin{enumerate}
                \item Fetch: Ambil instruksi dari memori
                \item Decode: Terjemahkan instruksi
                \item Execute: Jalankan instruksi
                \item Store: Simpan hasil (jika ada)
            \end{enumerate}
            
        \column{0.48\textwidth}
            \begin{center}
                \includegraphics[width=0.70\textwidth]{figure/fig8-cpuinscycle.png}
            \end{center}
    \end{columns}
\end{frame}

% --- Frame: Contoh Kasus Nyata ---
\begin{frame}{Contoh Kasus Nyata}
    \framesubtitle{Instruksi Penjumlahan Sederhana}
    
    \begin{center}
        Instruksi: \texttt{ADD A, B} (Tambahkan nilai B ke A)
    \end{center}
    
    \begin{enumerate}
        \item \textbf{Fetch}: CPU mengambil kode instruksi \texttt{ADD A, B} dari alamat 0x1000 di memori
        \item \textbf{Decode}: CPU mengidentifikasi bahwa ini adalah operasi penjumlahan dengan operand A dan B
        \item \textbf{Operand Fetch}: CPU mengambil nilai A (5) dari alamat 0x2000 dan B (3) dari alamat 0x2004
        \item \textbf{Execute}: CPU melakukan 5 + 3 = 8
        \item \textbf{Store}: CPU menyimpan hasil 8 kembali ke alamat memori A (0x2000)
    \end{enumerate}
\end{frame}



% --- Frame: Basic Hardware Considerations ---
\begin{frame}{Pertimbangan Hardware Dasar}
    \framesubtitle{Hierarki Akses Memori}
    \begin{itemize}
        \item \textbf{Akses Langsung CPU:} Terbatas pada register dan memori utama
            \begin{itemize}
                \item Tidak ada akses langsung ke disk atau penyimpanan sekunder lainnya
            \end{itemize}
        \item \textbf{Perbedaan Kecepatan Akses:}
            \begin{itemize}
                \item \textbf{Register:} Dapat diakses dalam satu siklus clock CPU
                \item \textbf{Memori Utama:} Membutuhkan beberapa siklus clock
            \end{itemize}
        \item \textbf{Solusi untuk Kesenjangan Performa:}
            \begin{itemize}
                \item Cache memori antara CPU dan memori utama
                \item Terletak pada chip CPU untuk akses cepat
                \item Dikelola oleh hardware (bukan OS)
            \end{itemize}
    \end{itemize}
\end{frame}

% --- Frame: Memory Protection ---
\begin{frame}{Proteksi Memori}
    \framesubtitle{Memastikan Operasi Sistem yang Benar}
    
    \begin{columns}[T]
        \column{0.48\textwidth}
            \textbf{Kebutuhan Proteksi:}
            \begin{itemize}
                \item Melindungi OS dari proses pengguna
                \item Melindungi proses pengguna satu sama lain
                \item Implementasi berbasis hardware untuk efisiensi
            \end{itemize}
            
            \textbf{Ruang Memori Per-Proses:}
            \begin{itemize}
                \item Fundamental untuk eksekusi konkuren
                \item Mendefinisikan rentang alamat yang legal
            \end{itemize}
            
        \column{0.48\textwidth}
            \begin{center}
                \includegraphics[width=0.75\textwidth]{figure/fig8-1.png} % Replace with actual base/limit register image
            \end{center}
    \end{columns}
\end{frame}

% --- Frame: Memory Protection Implementation ---
\begin{frame}{Implementasi Proteksi Memori}
    \framesubtitle{Register Base dan Limit}
    
    \begin{block}{Cara Kerjanya}
        \begin{itemize}
            \item \textbf{Register Base:} Menyimpan alamat memori fisik legal terkecil
            \item \textbf{Register Limit:} Menentukan ukuran rentang alamat
            \item Contoh: Base = 300040, Limit = 120900
            \item Alamat legal: 300040 sampai 420939
        \end{itemize}
    \end{block}
    
    \begin{alertblock}{Mekanisme Proteksi}
        \begin{itemize}
            \item Hardware CPU membandingkan setiap alamat yang dihasilkan dalam mode pengguna
            \item Upaya akses ilegal akan trap ke OS sebagai error fatal
            \item Mencegah modifikasi kode/data OS atau memori pengguna lain
            \item Register base dan limit hanya dapat diisi oleh OS menggunakan instruksi khusus
        \end{itemize}
    \end{alertblock}
\end{frame}

% --- Frame: Operating System Memory Access ---
\begin{frame}{Akses Memori Sistem Operasi}
    \framesubtitle{Operasi Istimewa}
    
    \begin{columns}[T]
        \column{0.48\textwidth}
            \textbf{Hak Istimewa OS:}
            \begin{itemize}
                \item Akses tak terbatas ke semua memori dalam mode kernel
                \item Dapat mengakses memori OS dan memori pengguna
            \end{itemize}
            
            \textbf{Operasi Memori Penting untuk OS:}
            \begin{itemize}
                \item Memuat program pengguna ke dalam memori
                \item Mencatat program jika terjadi error
                \item Mengakses dan memodifikasi parameter system call
            \end{itemize}
            
        \column{0.48\textwidth}
            \textbf{Operasi Memori untuk Multitasking:}
            \begin{itemize}
                \item Context switching antar proses
                \item Menyimpan status proses dari register ke memori
                \item Memuat konteks proses berikutnya dari memori ke register
                \item Melakukan I/O ke dan dari memori pengguna
                \item Menyediakan layanan sistem lainnya
            \end{itemize}
    \end{columns}
\end{frame}

\section{Address Binding}

% --- Frame: Introduction to Address Binding ---
\begin{frame}{Address Binding}
    \framesubtitle{Konsep Dasar}
    
    \begin{block}{Definisi}
        Address binding adalah proses mengaitkan alamat memori pada instruksi dan data program.
    \end{block}
    
    \begin{itemize}
        \item Program berada di disk sebagai berkas biner yang dapat dieksekusi
        \item Untuk dijalankan, program harus dimuat ke dalam memori
        \item Selama eksekusi, program mengakses instruksi dan data dari memori
        \item Setelah terminasi, memori yang digunakan program dibebaskan
        \item Sistem operasi memungkinkan proses ditempatkan pada berbagai bagian memori fisik
    \end{itemize}
    
    \begin{alertblock}{Fokus Utama}
        Memahami kapan dan bagaimana alamat memori diresolved selama siklus hidup program.
    \end{alertblock}
\end{frame}

% --- Frame: Address Spaces ---
\begin{frame}{Ruang Alamat dalam Eksekusi Program}
    \framesubtitle{Tahapan Transformasi Alamat}
    
    \begin{columns}[T]
        \column{0.60\textwidth}
            \begin{itemize}
                \item \textbf{Alamat Simbolik:}
                \begin{itemize}
                    \item Variabel dan label dalam kode sumber (seperti \texttt{count})
                \end{itemize}
                
                \item \textbf{Alamat Relatif (Relocatable):}
                \begin{itemize}
                    \item Alamat offset ("14 byte dari awal modul")
                    \item Dihasilkan oleh compiler
                \end{itemize}
                
                \item \textbf{Alamat Absolut (Fisik):}
                \begin{itemize}
                    \item Alamat memori fisik aktual (seperti 74014)
                    \item Alamat memori yang sebenarnya digunakan CPU
                \end{itemize}
            \end{itemize}
            
        \column{0.35\textwidth}
            \begin{center}
                \includegraphics[width=0.7\textwidth]{figure/fig8-3.png}
            \end{center}
    \end{columns}
\end{frame}

% --- Frame: Proses Kompilasi dan Eksekusi ---
\begin{frame}{Proses Kompilasi dan Eksekusi}
    \framesubtitle{Dari Kode Sumber hingga Program dalam Memori}
    
    \begin{block}{Tahapan Penting}
        \begin{itemize}
            \item Proses mengubah program dari kode sumber menjadi program yang berjalan di memori melibatkan beberapa tahap
            \item Setiap tahap melakukan transformasi alamat tertentu
            \item Proses ini menentukan bagaimana alamat pada kode sumber akhirnya diterjemahkan menjadi alamat memori fisik
        \end{itemize}
    \end{block}
\end{frame}

% --- Frame: Tahap 1 - Kompilasi ---
\begin{frame}{Tahap 1: Kompilasi}
    \framesubtitle{Dari Source Program ke Object File}
    
    \begin{itemize}
        \item \textbf{Input}: Source program (kode sumber)
        \begin{itemize}
            \item Berisi variabel dan fungsi dengan \textbf{alamat simbolik}
            \item Contoh: \texttt{int count;}, \texttt{void calculate()}
        \end{itemize}
        
        \item \textbf{Proses}: Compiler
        \begin{itemize}
            \item Mengubah kode sumber menjadi kode assembly
            \item Mengganti alamat simbolik dengan \textbf{alamat relatif}
        \end{itemize}
        
        \item \textbf{Output}: Object file
        \begin{itemize}
            \item Kode mesin dengan alamat relatif
            \item Format: \texttt{.o} (Unix/Linux), \texttt{.obj} (Windows)
        \end{itemize}
    \end{itemize}
\end{frame}

% --- Frame: Tahap 2 - Linking ---
\begin{frame}{Tahap 2: Linking}
    \framesubtitle{Dari Object File ke Executable File}
    
    \begin{itemize}
        \item \textbf{Input}: Multiple object files dan library files
        \begin{itemize}
            \item Object file dari kode sumber utama
            \item Object file dari modul tambahan
            \item Library standar dan khusus
        \end{itemize}
        
        \item \textbf{Proses}: Linker
        \begin{itemize}
            \item Menggabungkan semua object file
            \item Me-resolve referensi eksternal
            \item Menghasilkan ruang alamat tunggal dan terstruktur
        \end{itemize}
        
        \item \textbf{Output}: Executable file
        \begin{itemize}
            \item Program yang siap dieksekusi
            \item Format: \texttt{ELF} (Linux), \texttt{.exe} (Windows)
        \end{itemize}
    \end{itemize}
\end{frame}

% --- Frame: Tahap 3 - Loading ---
\begin{frame}{Tahap 3: Loading}
    \framesubtitle{Dari Executable File ke Program dalam Memori}
    
    \begin{itemize}
        \item \textbf{Input}: Executable file
        \begin{itemize}
            \item Program yang siap dieksekusi dari disk
        \end{itemize}
        
        \item \textbf{Proses}: Loader
        \begin{itemize}
            \item Membaca executable file ke memori
            \item Mengalokasikan ruang memori untuk program
            \item Mengubah alamat relatif menjadi \textbf{alamat absolut}
        \end{itemize}
        
        \item \textbf{Tambahan}: Dynamic Linking
        \begin{itemize}
            \item Menghubungkan \textit{dynamically linked libraries} (DLL)
            \item Contoh: \texttt{.so} (Linux), \texttt{.dll} (Windows)
            \item Menghemat memori dengan berbagi kode
        \end{itemize}
    \end{itemize}
\end{frame}

% --- Frame: Program dalam Memori ---
\begin{frame}{Program dalam Memori}
    \framesubtitle{Hasil Akhir Proses}
    
    \textbf{Program yang Siap Dieksekusi:}
    \begin{itemize}
        \item Seluruh alamat kini adalah \textbf{alamat absolut}
        \item Program memiliki ruang alamat tertentu dalam memori fisik
        \item CPU dapat mengakses instruksi dan data langsung
        \item Sistem operasi mengatur eksekusi program
    \end{itemize}
    
    \textbf{Komponen dalam Memori:}
    \begin{itemize}
        \item Segmen kode (code segment)
        \item Segmen data (data segment)
        \item Stack untuk fungsi dan variabel lokal
        \item Heap untuk alokasi memori dinamis
    \end{itemize}
\end{frame}

% --- Frame: Ringkasan Proses ---
\begin{frame}{Ringkasan Proses Transformasi Program}
    \framesubtitle{Dari Kode Sumber hingga Eksekusi}
    
    \begin{enumerate}
        \item \textbf{Source Program}
        \begin{itemize}
            \item Kode yang ditulis dalam bahasa pemrograman
            \item Menggunakan \textbf{alamat simbolik}
        \end{itemize}
        
        \item \textbf{Compiler} $\Rightarrow$ \textbf{Object File}
        \begin{itemize}
            \item Mengubah alamat simbolik menjadi \textbf{alamat relatif}
        \end{itemize}
        
        \item \textbf{Linker} $\Rightarrow$ \textbf{Executable File}
        \begin{itemize}
            \item Menggabungkan object files dan library
            \item Menghasilkan program yang siap dieksekusi
        \end{itemize}
        
        \item \textbf{Loader + Dynamic Libraries} $\Rightarrow$ \textbf{Program in Memory}
        \begin{itemize}
            \item Mengubah alamat relatif menjadi \textbf{alamat absolut}
            \item Program siap dieksekusi oleh CPU
        \end{itemize}
    \end{enumerate}
\end{frame}

\section{Logical vs Physical Address Space}
% --- Frame: Logical vs Physical Address Spaces ---
\begin{frame}{Logical vs Physical Address Space}
    \framesubtitle{Konsep Dasar}
    
    \begin{columns}[T]
        \column{0.48\textwidth}
            \textbf{Logical Address (Virtual Address):}
            \begin{itemize}
                \item Alamat yang dihasilkan oleh CPU saat eksekusi
                \item Referensi ke lokasi memori independen dari alokasi memori sebenarnya
                \item Pengguna program hanya melihat ruang alamat logis
                \item Juga disebut sebagai alamat virtual
            \end{itemize}
            
        \column{0.48\textwidth}
            \textbf{Physical Address:}
            \begin{itemize}
                \item Alamat lokasi sebenarnya di memori fisik
                \item Alamat yang benar-benar diakses pada unit memori
                \item Tidak terlihat oleh program pengguna
                \item Hanya digunakan oleh unit memori untuk akses aktual
            \end{itemize}
    \end{columns}
\end{frame}

% --- Frame: Memory Management Unit (MMU) ---
\begin{frame}{Memory Management Unit (MMU)}
    \framesubtitle{Perangkat Hardware untuk Translasi Alamat}
    
    \begin{block}{Definisi}
        MMU adalah komponen hardware yang menerjemahkan alamat logis menjadi alamat fisik secara real-time saat eksekusi program.
    \end{block}
    
    \begin{columns}[T]
        \column{0.58\textwidth}
            \textbf{Fungsi Utama MMU:}
            \begin{itemize}
                \item Mentranslasi setiap alamat logis yang dihasilkan CPU
                \item Menyembunyikan detail memori fisik dari program
                \item Menyediakan mekanisme proteksi memori
                \item Mendukung berbagai skema manajemen memori
            \end{itemize}
            
        \column{0.38\textwidth}
            \begin{center}
            \textbf{MMU}
            \begin{itemize}
                \item CPU $\rightarrow$ MMU $\rightarrow$ Memory
                \item Logical Address $\rightarrow$ Physical Address
            \end{itemize}
            \end{center}
    \end{columns}
\end{frame}

% --- Frame: Logical vs Physical Address - Daily Life Illustration ---
\begin{frame}{Ilustrasi Sehari-hari: Logical vs Physical Address}
    \framesubtitle{Perpustakaan Kota}
    
    \begin{columns}[T]
        \column{0.48\textwidth}
            \textbf{Logical Address:}
            \begin{itemize}
                \item \textbf{Katalog perpustakaan} yang mencantumkan buku-buku
                \item Pengunjung mencari buku berdasarkan \textbf{sistem katalog}
                \item Pengunjung tidak perlu tahu di rak mana fisik buku berada
                \item Katalog bisa tetap sama meskipun tata letak perpustakaan berubah
            \end{itemize}
            
        \column{0.48\textwidth}
            \textbf{Physical Address:}
            \begin{itemize}
                \item \textbf{Lokasi sebenarnya} buku pada rak tertentu
                \item Contoh: "Lantai 2, Rak E, Baris 3"
                \item Pustakawan menggunakan pemetaan dari nomor katalog ke lokasi fisik
                \item Lokasi fisik dapat berubah (reorganisasi) tanpa mengubah katalog
            \end{itemize}
    \end{columns}
    
    \vspace{0.3cm}
    \begin{alertblock}{Analogi MMU}
        Pustakawan bertindak seperti MMU - menerjemahkan referensi katalog (alamat logis) menjadi lokasi buku fisik (alamat fisik) saat pengunjung meminta buku.
    \end{alertblock}
\end{frame}

% --- Frame: Relocation Register ---
\begin{frame}{Relocation Register}
    \framesubtitle{Mekanisme Sederhana untuk Dynamic Relocation}
    
    \begin{columns}[T]
        \column{0.48\textwidth}
            \textbf{Definisi:}
            \begin{itemize}
                \item Register perangkat keras yang menyimpan nilai \textbf{base address}
                \item Alamat awal dari area memori yang dialokasikan untuk sebuah proses
                \item Digunakan untuk menghitung alamat fisik dari alamat logis
            \end{itemize}
            
            \textbf{Cara Kerja Sederhana:}
            \begin{itemize}
                \item \texttt{Alamat Fisik = Relocation Register + Alamat Logis}
            \end{itemize}
            
        \column{0.48\textwidth}
            \begin{center}
                \textbf{Prinsip Relocation Register}
                \vspace{0.5cm}
                \begin{tabular}{|c|c|c|}
                    \hline
                    \textbf{Alamat Logis} & \textbf{+} & \textbf{Relocation Reg.} \\
                    \hline
                    346 & + & 14000 \\
                    \hline
                    \multicolumn{3}{|c|}{\textbf{= 14346 (Alamat Fisik)}} \\
                    \hline
                \end{tabular}
            \end{center}
    \end{columns}
\end{frame}

% --- Frame: Dynamic Relocation Scheme ---
\begin{frame}{Dynamic Relocation Scheme}
    \framesubtitle{Menggunakan Relocation Register}
    
    \begin{columns}[T]
        \column{0.55\textwidth}
            \textbf{Prinsip Kerja:}
                        \begin{itemize}
                            \item Setiap alamat logis program ditambahkan ke nilai relocation register
                            \item Hardware menangani penambahan secara otomatis saat runtime
                        \end{itemize}
                        
                        \textbf{Manfaat:}
                        \begin{itemize}
                            \item Program dapat dimuat di lokasi memori berbeda tanpa modifikasi kode
                            \item Relokasi dilakukan secara transparan
                        \end{itemize}
            
        \column{0.42\textwidth}
            \begin{block}{Contoh Numerik}
                \begin{itemize}
                    \item Relocation register = 14000
                    \item Alamat logis = 346
                    \item Alamat fisik = 14000 + 346 = 14346
                \end{itemize}
            \end{block}
            
            \begin{alertblock}{Keamanan}
                Kombinasi dengan register limit mencegah proses mengakses di luar area yang dialokasikan.
            \end{alertblock}
    \end{columns}
\end{frame}

% --- Frame: Demonstration of Relocation Register ---
\begin{frame}{Demonstrasi Relocation Register}
    \framesubtitle{Kasus Nyata: Multitasking OS}
    
    \begin{block}{Skenario}
        Sistem operasi memuat tiga program berbeda ke dalam memori secara bersamaan.
    \end{block}

    \begin{center}
        \begin{tabular}{|l|c|c|c|}
            \hline
            \textbf{Program} & \textbf{Base Address} & \textbf{Limit} & \textbf{Range di Memori Fisik} \\
            \hline
            Word Processor & 20000 & 8000 & 20000-27999 \\
            \hline
            Game & 35000 & 12000 & 35000-46999 \\
            \hline
            Browser & 55000 & 15000 & 55000-69999 \\
            \hline
        \end{tabular}
    \end{center}
\end{frame}

% --- Frame: Praktik Relocation Register - Program View ---
\begin{frame}{Relocation Register dalam Praktik}
    \framesubtitle{Perspektif Program}
    
    \begin{columns}[T]
        \column{0.48\textwidth}
            \textbf{Kode Game - Perspektif Program:}
            \begin{itemize}
                \item Program ditulis dan dikompilasi dengan asumsi dimulai dari alamat 0
                \item Code segment: 0-4999
                \item Data segment: 5000-9999
                \item Stack: 10000-11999
            \end{itemize}
            \begin{block}{Instruksi dalam Program}
                LOAD dari alamat 6500 \\
                JUMP ke alamat 2500
            \end{block}
            
        \column{0.48\textwidth}
            \textbf{Relocation dalam Hardware:}
            \begin{itemize}
                \item Relocation register = 35000
                \item LOAD dari 35000+6500 = 41500
                \item JUMP ke 35000+2500 = 37500
            \end{itemize}
            \begin{alertblock}{Keuntungan}
                Program tidak tahu alamat fisiknya. Bisa dipindahkan ke lokasi memori berbeda tanpa perubahan kode.
            \end{alertblock}
    \end{columns}
\end{frame}


% --- Frame: Perlindungan Memori dengan Base dan Limit ---
\begin{frame}{Perlindungan Memori dengan Base dan Limit}
    \framesubtitle{Contoh Kasus Operasional}
    
    \textbf{Skenario:} Browser mencoba mengakses alamat yang tidak valid
    
    \begin{columns}[T]
        \column{0.48\textwidth}
            \textbf{Kasus 1: Akses Valid}
            \begin{itemize}
                \item Alamat logis: 12500
                \item Limit: 15000
                \item Akses valid karena 12500 < 15000
                \item Alamat fisik: 55000 (base) + 12500 = 67500
                \item Operasi berhasil
            \end{itemize}
            
        \column{0.48\textwidth}
            \textbf{Kasus 2: Akses Tidak Valid}
            \begin{itemize}
                \item Alamat logis: 16200
                \item Limit: 15000
                \item Akses tidak valid karena 16200 > 15000
                \item Hardware mendeteksi pelanggaran batas
                \item Menghasilkan segmentation fault
                \item OS menghentikan program yang melanggar
            \end{itemize}
    \end{columns}
\end{frame}

% --- Frame: Visualisasi Dynamic Relocation ---
\begin{frame}{Visualisasi Dynamic Relocation}
    \framesubtitle{Pemetaan Ruang Alamat dengan Relocation Register}
    
    \begin{center}
        \begin{tabular}{|c|c|}
            \hline
            \textbf{Logical Address Space} & \textbf{Physical Address Space} \\
            \hline
            0 & 14000 (Base) \\
            \hline
            100 & 14100 \\
            \hline
            200 & 14200 \\
            \hline
            ... & ... \\
            \hline
            4000 (Limit) & 18000 \\
            \hline
        \end{tabular}
    \end{center}
    
    \begin{itemize}
        \item Program ditulis seolah-olah dimuat pada alamat 0
        \item Relocation register menambahkan offset ke setiap alamat
        \item Program dapat dieksekusi di lokasi memori manapun tanpa modifikasi
        \item Berfungsi seperti "memindahkan" seluruh ruang alamat logis ke lokasi fisik
    \end{itemize}
\end{frame}

\section{Memory Allocation}

% --- Frame: Memory Allocation Basics ---
\begin{frame}{Alokasi Memori}
    \framesubtitle{Konsep Dasar}
    
    \begin{block}{Definisi}
        Alokasi memori adalah proses pemberian ruang memori kepada proses yang membutuhkan sumber daya memori untuk eksekusi.
    \end{block}
    
    \begin{itemize}
        \item Sistem operasi melakukan pembagian memori ke partisi dengan ukuran bervariasi
        \item Setiap partisi dapat berisi tepat satu proses
        \item OS menyimpan tabel yang menunjukkan bagian memori yang tersedia dan yang sudah terpakai
        \item Awalnya, semua memori tersedia sebagai satu blok besar (disebut \textit{hole})
        \item Seiring waktu, memori akan berisi kumpulan \textit{holes} dengan berbagai ukuran
    \end{itemize}
\end{frame}

% --- Frame: Memory Allocation Process ---
\begin{frame}{Proses Alokasi Memori}
    \framesubtitle{Mekanisme Dasar}
    
    \begin{columns}[T]
        \column{0.48\textwidth}
            \textbf{Proses Alokasi:}
            \begin{itemize}
                \item OS mempertimbangkan kebutuhan memori proses dan jumlah memori yang tersedia
                \item Ketika proses dialokasikan ruang, proses dimuat ke memori
                \item Setelah termuat, proses dapat berkompetisi untuk waktu CPU
                \item Ketika proses berakhir, memori dibebaskan kembali
            \end{itemize}
            
        \column{0.48\textwidth}
            \textbf{Ketika Memori Tidak Cukup:}
            \begin{itemize}
                \item Opsi 1: Tolak proses dan berikan pesan error
                \item Opsi 2: Tempatkan proses ke dalam \textit{wait queue}
                \item Ketika memori dibebaskan, OS memeriksa \textit{wait queue}
                \item Jika memori cukup, proses yang menunggu dapat dilayani
            \end{itemize}
    \end{columns}
\end{frame}

% --- Frame: Dynamic Storage Allocation Problem ---
\begin{frame}{Dynamic Storage Allocation Problem}
    \framesubtitle{Tantangan Alokasi Memori}
    
    \begin{block}{Proses Umum Alokasi}
        \begin{itemize}
            \item Ketika proses tiba dan membutuhkan memori, sistem mencari \textit{hole} yang cukup besar
            \item Jika \textit{hole} terlalu besar, dipecah menjadi dua bagian:
            \begin{itemize}
                \item Satu bagian untuk proses yang datang
                \item Bagian lain dikembalikan ke kumpulan \textit{holes}
            \end{itemize}
            \item Ketika proses berakhir, blok memori dikembalikan sebagai \textit{hole}
            \item Jika \textit{hole} baru berdekatan dengan \textit{holes} lain, keduanya digabungkan
        \end{itemize}
    \end{block}
    
    \begin{alertblock}{Masalah Inti}
        Bagaimana memenuhi permintaan memori dengan ukuran n dari daftar \textit{free holes} yang tersedia?
    \end{alertblock}
\end{frame}

% --- Frame: Memory Allocation Strategies ---
\begin{frame}{Strategi Alokasi Memori}
    \framesubtitle{Pendekatan Umum}
    
    \begin{block}{Strategi-strategi Utama}
        Ada tiga strategi umum yang digunakan untuk memilih \textit{free hole} dari kumpulan \textit{holes} yang tersedia:
        \begin{enumerate}
            \item \textbf{First-fit}
            \item \textbf{Best-fit}
            \item \textbf{Worst-fit}
        \end{enumerate}
    \end{block}
    
    \begin{alertblock}{Tujuan}
        Setiap strategi berusaha mengelola fragmentasi memori dan meningkatkan efisiensi penggunaan memori dengan cara yang berbeda.
    \end{alertblock}
\end{frame}

% --- Frame: First-Fit Strategy ---
\begin{frame}{First-Fit Strategy}
    \framesubtitle{Strategi Alokasi Pertama-Cocok}
    
    \begin{block}{Definisi}
        \textit{First-fit} mengalokasikan \textit{hole} pertama yang ditemukan yang cukup besar untuk menampung proses.
    \end{block}
    
    \begin{columns}[T]
        \column{0.48\textwidth}
            \textbf{Cara Kerja:}
            \begin{itemize}
                \item Memindai daftar memori dari awal
                \item Memilih \textit{hole} pertama yang ukurannya ≥ kebutuhan proses
                \item Algoritma berhenti segera setelah menemukan \textit{hole} yang cocok
            \end{itemize}
            
        \column{0.48\textwidth}
            \textbf{Karakteristik:}
            \begin{itemize}
                \item \textbf{Kecepatan:} Sangat cepat karena pencarian minimal
                \item \textbf{Implementasi:} Sederhana
                \item \textbf{Efisiensi:} Terkadang meninggalkan \textit{holes} kecil di memori awal
                \item \textbf{Penggunaan:} Sering digunakan dalam praktik
            \end{itemize}
    \end{columns}
\end{frame}

% --- Frame: Best-Fit Strategy ---
\begin{frame}{Best-Fit Strategy}
    \framesubtitle{Strategi Alokasi Terbaik-Cocok}
    
    \begin{block}{Definisi}
        \textit{Best-fit} mengalokasikan \textit{hole} terkecil yang masih cukup besar untuk menampung proses.
    \end{block}
    
    \begin{columns}[T]
        \column{0.48\textwidth}
            \textbf{Cara Kerja:}
            \begin{itemize}
                \item Memindai seluruh daftar memori
                \item Mencari \textit{hole} dengan ukuran paling mendekati kebutuhan proses
                \item Memilih \textit{hole} dengan sisa ruang terkecil
            \end{itemize}
            
        \column{0.48\textwidth}
            \textbf{Karakteristik:}
            \begin{itemize}
                \item \textbf{Kecepatan:} Lebih lambat (harus memindai seluruh daftar)
                \item \textbf{Fragmentasi:} Meminimalkan fragmentasi eksternal
                \item \textbf{Efisiensi:} Sering menghasilkan \textit{holes} yang terlalu kecil untuk digunakan
                \item \textbf{Pencarian:} Lebih intensif dibanding \textit{first-fit}
            \end{itemize}
    \end{columns}
\end{frame}

% --- Frame: Worst-Fit Strategy ---
\begin{frame}{Worst-Fit Strategy}
    \framesubtitle{Strategi Alokasi Terburuk-Cocok}
    
    \begin{block}{Definisi}
        \textit{Worst-fit} mengalokasikan \textit{hole} terbesar yang tersedia untuk menampung proses.
    \end{block}
    
    \begin{columns}[T]
        \column{0.48\textwidth}
            \textbf{Cara Kerja:}
            \begin{itemize}
                \item Memindai seluruh daftar memori
                \item Mencari \textit{hole} dengan ukuran paling besar
                \item Memilih \textit{hole} dengan sisa ruang terbesar
            \end{itemize}
            
        \column{0.48\textwidth}
            \textbf{Karakteristik:}
            \begin{itemize}
                \item \textbf{Kecepatan:} Lambat (harus memindai seluruh daftar)
                \item \textbf{Pemikiran:} Meninggalkan \textit{holes} yang cukup besar untuk digunakan
                \item \textbf{Realitas:} Cepat menghabiskan \textit{holes} besar
                \item \textbf{Performa:} Secara umum lebih buruk dari \textit{first-fit} dan \textit{best-fit}
            \end{itemize}
    \end{columns}
\end{frame}

% --- Frame: Perbandingan Strategi ---
\begin{frame}{Perbandingan Strategi Alokasi}
    \framesubtitle{First-fit vs Best-fit vs Worst-fit}
    
    \begin{center}
        \begin{tabular}{|l|c|c|c|}
            \hline
            \textbf{Aspek} & \textbf{First-fit} & \textbf{Best-fit} & \textbf{Worst-fit} \\
            \hline
            Kecepatan & Tinggi & Sedang & Rendah \\
            \hline
            Fragmentasi & Sedang & Tinggi (banyak \textit{tiny holes}) & Rendah \\
            \hline
            Pemanfaatan memori & Baik & Baik & Kurang baik \\
            \hline
            Kompleksitas implementasi & Rendah & Sedang & Sedang \\
            \hline
            Performa keseluruhan & Sangat baik & Baik & Kurang baik \\
            \hline
        \end{tabular}
    \end{center}
    
    \begin{alertblock}{Catatan Penting}
        Studi simulasi menunjukkan bahwa \textit{first-fit} dan \textit{best-fit} lebih baik daripada \textit{worst-fit} dalam hal pemanfaatan memori dan waktu.
    \end{alertblock}
\end{frame}

% --- Frame: Contoh Kasus - Strategi Alokasi Memori ---
\begin{frame}{Contoh Kasus: Strategi Alokasi Memori}
    \framesubtitle{Perbandingan Algoritma dalam Situasi Nyata}
    
    \begin{block}{Kasus}
        Diberikan partisi memori sebagai berikut:
        \begin{itemize}
            \item 100K, 500K, 200K, 300K, dan 600K (berurutan)
        \end{itemize}
        
        Bagaimana algoritma First-fit, Best-fit, dan Worst-fit akan menempatkan proses berikut?
        \begin{itemize}
            \item 212K, 417K, 112K, dan 624K (berurutan)
        \end{itemize}
    \end{block}
    
    \begin{alertblock}{Pertanyaan}
        Algoritma mana yang membuat penggunaan memori paling efisien?
    \end{alertblock}
\end{frame}

% --- Frame: First-Fit Approach Solution ---
\begin{frame}{Solusi First-Fit Approach}
    \framesubtitle{Proses Alokasi Langkah demi Langkah}
    
    \begin{columns}[T]
        \column{0.48\textwidth}
            \textbf{Awal:} 100K, 500K, 200K, 300K, 600K
            
            \textbf{Proses 212K:}
            \begin{itemize}
                \item Tidak muat di 100K
                \item Muat di 500K $\rightarrow$ Tersisa 288K
                \item Status: [100K][288K][200K][300K][600K]
            \end{itemize}
            
            \textbf{Proses 417K:}
            \begin{itemize}
                \item Tidak muat di 100K, 288K
                \item Tidak muat di 200K
                \item Muat di 300K $\rightarrow$ Tidak ada sisa
                \item Status: [100K][288K][200K][0][600K]
            \end{itemize}
            
        \column{0.48\textwidth}
            \textbf{Proses 112K:}
            \begin{itemize}
                \item Muat di 100K $\rightarrow$ Tidak ada sisa
                \item Status: [0][288K][200K][0][600K]
            \end{itemize}
            
            \textbf{Proses 624K:}
            \begin{itemize}
                \item Tidak muat di 288K, 200K
                \item Muat di 600K $\rightarrow$ Tidak ada sisa
                \item Status: [0][288K][200K][0][0]
            \end{itemize}
            
            \textbf{Akhir:} Memori yang tersisa: 288K + 200K = 488K
    \end{columns}
\end{frame}

% --- Frame: Best-Fit Approach Solution ---
\begin{frame}{Solusi Best-Fit Approach}
    \framesubtitle{Proses Alokasi Langkah demi Langkah}
    
    \begin{columns}[T]
        \column{0.48\textwidth}
            \textbf{Awal:} 100K, 500K, 200K, 300K, 600K
            
            \textbf{Proses 212K:}
            \begin{itemize}
                \item Best-fit untuk 212K adalah 300K 
                \item Tersisa 88K di partisi tersebut
                \item Status: [100K][500K][200K][88K][600K]
            \end{itemize}
            
            \textbf{Proses 417K:}
            \begin{itemize}
                \item Best-fit untuk 417K adalah 500K
                \item Tersisa 83K di partisi tersebut
                \item Status: [100K][83K][200K][88K][600K]
            \end{itemize}
            
        \column{0.48\textwidth}
            \textbf{Proses 112K:}
            \begin{itemize}
                \item Best-fit untuk 112K adalah 200K
                \item Tersisa 88K di partisi tersebut
                \item Status: [100K][83K][88K][88K][600K]
            \end{itemize}
            
            \textbf{Proses 624K:}
            \begin{itemize}
                \item Best-fit untuk 624K adalah 600K
                \item Tidak ada sisa memori
                \item Status: [100K][83K][88K][88K][0]
            \end{itemize}
            
            \textbf{Akhir:} Memori yang tersisa: 100K + 83K + 88K + 88K = 359K
    \end{columns}
\end{frame}

% --- Frame: Worst-Fit Approach Solution ---
\begin{frame}{Solusi Worst-Fit Approach}
    \framesubtitle{Proses Alokasi Langkah demi Langkah}
    
    \begin{columns}[T]
        \column{0.48\textwidth}
            \textbf{Awal:} 100K, 500K, 200K, 300K, 600K
            
            \textbf{Proses 212K:}
            \begin{itemize}
                \item Worst-fit untuk 212K adalah 600K
                \item Tersisa 388K di partisi tersebut
                \item Status: [100K][500K][200K][300K][388K]
            \end{itemize}
            
            \textbf{Proses 417K:}
            \begin{itemize}
                \item Worst-fit untuk 417K adalah 500K
                \item Tersisa 83K di partisi tersebut
                \item Status: [100K][83K][200K][300K][388K]
            \end{itemize}
            
        \column{0.48\textwidth}
            \textbf{Proses 112K:}
            \begin{itemize}
                \item Worst-fit untuk 112K adalah 388K
                \item Tersisa 276K di partisi tersebut
                \item Status: [100K][83K][200K][300K][276K]
            \end{itemize}
            
            \textbf{Proses 624K:}
            \begin{itemize}
                \item Tidak ada partisi yang cukup besar
                \item Proses harus menunggu
                \item Status: [100K][83K][200K][300K][276K]
            \end{itemize}
            
            \textbf{Akhir:} Memori yang tersisa: 100K + 83K + 200K + 300K + 276K = 959K
    \end{columns}
\end{frame}

% --- Frame: Perbandingan dan Kesimpulan ---
\begin{frame}{Perbandingan dan Kesimpulan}
    \framesubtitle{Analisis Hasil Algoritma}
    
    \begin{columns}[T]
        \column{0.48\textwidth}
            \textbf{Ringkasan Hasil:}
                
                \begin{center}
                    \begin{tabular}{|l|c|c|}
                      \hline
                      \textbf{Algoritma} & \textbf{Proses} & \textbf{Memori} \\
                      & \textbf{Berhasil} & \textbf{Tersisa} \\
                      \hline
                      First-fit & 4 & 488K \\
                      \hline
                      Best-fit & 4 & 359K \\
                      \hline
                      Worst-fit & 3 & 959K \\
                      \hline
                    \end{tabular}
                \end{center}
            
        \column{0.48\textwidth}
            \textbf{Kesimpulan:}
            \begin{itemize}
                \item Best-fit menghasilkan penggunaan memori yang paling efisien dengan memori tersisa paling sedikit
                \item First-fit menggunakan memori dengan efisiensi sedang
                \item Worst-fit tidak dapat mengakomodasi semua proses dan memiliki fragmentasi tertinggi
            \end{itemize}
            
            \begin{alertblock}{Kunci Takeaway}
                Untuk kasus ini, Best-fit adalah yang paling efisien dalam penggunaan memori, sementara Worst-fit adalah yang paling tidak efisien.
            \end{alertblock}
    \end{columns}
\end{frame}

% --- Frame: Visualisasi Perbandingan ---
\begin{frame}{Visualisasi Perbandingan}
    \framesubtitle{Pemanfaatan Memori pada Setiap Algoritma}
        
        \begin{columns}[T]
            \column{0.48\textwidth}
                \begin{block}{First-fit}
                    \begin{itemize}
                        \item Memori awal: 1700K
                        \item Memori terpakai: 1365K
                        \item Efisiensi: 80,3\%
                    \end{itemize}
                \end{block}
                
                \begin{block}{Best-fit}
                    \begin{itemize}
                        \item Memori awal: 1700K
                        \item Memori terpakai: 1365K
                        \item Efisiensi: 80,3\%
                    \end{itemize}
                \end{block}
                
            \column{0.48\textwidth}
                \begin{block}{Worst-fit}
                    \begin{itemize}
                        \item Memori awal: 1700K
                        \item Memori terpakai: 741K
                        \item Efisiensi: 43,6\%
                    \end{itemize}
                \end{block}
                
                \begin{alertblock}{Kesimpulan}
                    \begin{itemize}
                        \item First-fit dan Best-fit memiliki efisiensi yang sama (80,3\%)
                        \item Worst-fit memiliki efisiensi paling rendah (43,6\%)
                        \item First-fit lebih mudah diimplementasikan dan lebih cepat
                    \end{itemize}
                \end{alertblock}
        \end{columns}
\end{frame}

\section{Fragmentasi Memori}

% --- Frame: Pengenalan Fragmentasi ---
\begin{frame}{Fragmentasi Memori}
    \framesubtitle{Konsekuensi dari Alokasi Dinamis}
    
    \begin{block}{Definisi}
        Fragmentasi adalah kondisi di mana memori terbagi menjadi banyak potongan kecil yang sulit dimanfaatkan secara efisien.
    \end{block}
    
    \begin{itemize}
        \item Muncul sebagai \textbf{efek samping} dari proses alokasi dan dealokasi memori
        \item Menyebabkan \textbf{pemborosan} memori meskipun memori total mencukupi
        \item Mengurangi efisiensi penggunaan sumber daya sistem
        \item Menjadi masalah serius dalam sistem yang berjalan dalam waktu lama
    \end{itemize}
\end{frame}

% --- Frame: Jenis Fragmentasi ---
\begin{frame}{Jenis Fragmentasi}
    \framesubtitle{External vs Internal Fragmentation}
    
    \begin{columns}[T]
        \column{0.48\textwidth}
            \textbf{External Fragmentation:}
            \begin{itemize}
                \item Total memori bebas cukup untuk permintaan, tetapi tidak bersebelahan
                \item Memori terpecah menjadi banyak \textit{holes} kecil
                \item Terjadi di antara blok-blok memori yang terpakai
                \item Contoh: Total 200KB bebas, tetapi terpecah menjadi blok 50KB, 75KB, dan 75KB
            \end{itemize}
            
        \column{0.48\textwidth}
            \textbf{Internal Fragmentation:}
            \begin{itemize}
                \item Memori yang dialokasikan lebih besar dari yang diminta
                \item Ruang ekstra tidak dapat digunakan oleh proses lain
                \item Terjadi di dalam blok memori terpakai
                \item Contoh: Proses meminta 90KB, sistem mengalokasikan 100KB, 10KB terbuang
            \end{itemize}
    \end{columns}
\end{frame}

% --- Frame: Visualisasi Fragmentasi Eksternal ---
\begin{frame}{Visualisasi External Fragmentation}
    \framesubtitle{Masalah dalam Alokasi Dinamis}
    
    \begin{center}
        \begin{tabular}{|c|c|c|c|c|c|c|c|c|}
        \hline
        \multicolumn{9}{|c|}{Memori} \\
        \hline
        Proses A & \cellcolor{gray!30}Hole & Proses B & \cellcolor{gray!30}Hole & Proses C & \cellcolor{gray!30}Hole & Proses D & \cellcolor{gray!30}Hole & Proses E \\
        \hline
        50K & 20K & 100K & 10K & 75K & 15K & 40K & 25K & 150K \\
        \hline
        \end{tabular}
    \end{center}
    
    \begin{alertblock}{Masalah}
        \begin{itemize}
            \item Total memori bebas: 20K + 10K + 15K + 25K = 70K
            \item Memori cukup untuk proses baru yang membutuhkan 65K
            \item Tetapi tidak ada \textit{hole} tunggal yang cukup besar
            \item Proses baru harus menunggu meskipun memori total mencukupi
        \end{itemize}
    \end{alertblock}
\end{frame}

% --- Frame: Visualisasi Fragmentasi Internal ---
\begin{frame}{Visualisasi Internal Fragmentation}
    \framesubtitle{Pemborosan dalam Alokasi}
    
    \begin{columns}[T]
        \column{0.48\textwidth}
            \begin{center}
                \begin{tabular}{|c|c|}
                \hline
                \multicolumn{2}{|c|}{Blok Memori 100K} \\
                \hline
                \cellcolor{blue!20}Proses A (92K) & \cellcolor{red!20}8K \\
                \hline
                \end{tabular}
                
                \vspace{0.5cm}
                
                \begin{tabular}{|c|c|}
                \hline
                \multicolumn{2}{|c|}{Blok Memori 64K} \\
                \hline
                \cellcolor{blue!20}Proses B (60K) & \cellcolor{red!20}4K \\
                \hline
                \end{tabular}
            \end{center}
            
        \column{0.48\textwidth}
            \begin{alertblock}{Masalah}
                \begin{itemize}
                    \item 8K + 4K = 12K memori terbuang
                    \item Tidak dapat digunakan oleh proses lain
                    \item Terjadi karena alokasi blok berukuran tetap
                    \item Sering terjadi pada sistem dengan alokasi berbasis unit atau halaman
                \end{itemize}
            \end{alertblock}
            
            \begin{block}{Penyebab Umum}
                \begin{itemize}
                    \item Kebutuhan alignment memori
                    \item Alokasi berbasis blok
                    \item Kebijakan alokasi yang tidak fleksibel
                \end{itemize}
            \end{block}
    \end{columns}
\end{frame}

% --- Frame: 50 Percent Rule ---
\begin{frame}{50 Percent Rule}
    \framesubtitle{Fenomena dalam External Fragmentation}
    
    \begin{block}{Definisi}
        Untuk alokasi dan dealokasi first-fit acak, hampir 1/3 blok memori akan terbuang karena fragmentasi.
    \end{block}
    
    \begin{columns}[T]
        \column{0.48\textwidth}
            \textbf{Formula:}
            \begin{itemize}
                \item Jika N blok terpakai, perkiraan jumlah \textit{holes} = (N+1)/2
                \item Jika memori rata-rata per proses = P, maka memori yang hilang karena fragmentasi ≈ P/2
                \item Artinya, fragmentasi eksternal = 1/3 memori total
            \end{itemize}
            
        \column{0.48\textwidth}
            \textbf{Implikasi:}
            \begin{itemize}
                \item Alokasi memori dinamis tanpa kompaksi akan kehilangan sekitar 1/3 memori
                \item Sistem harus mengalokasikan memori 50\% lebih besar dari kebutuhan sebenarnya
                \item Masalah menjadi lebih parah dalam sistem yang berjalan lama
                \item Performa sistem menurun seiring waktu
            \end{itemize}
    \end{columns}
\end{frame}

% --- Frame: Metode Mengatasi Fragmentasi ---
\begin{frame}{Mengatasi Fragmentasi}
    \framesubtitle{Teknik untuk Mengoptimalkan Penggunaan Memori}
    
    \begin{columns}[T]
        \column{0.48\textwidth}
            \textbf{Kompaksi (Compaction):}
            \begin{itemize}
                \item Menggeser memori untuk menggabungkan \textit{holes} kecil
                \item Membutuhkan kemampuan relokasi dinamis
                \item Mahal dalam hal overhead komputasi
                \item Contoh: Defragmentasi disk
            \end{itemize}
            
            \textbf{Segmentasi:}
            \begin{itemize}
                \item Membagi memori logis program menjadi segmen
                \item Segmen dapat dialokasikan secara terpisah
                \item Meningkatkan fleksibilitas alokasi
            \end{itemize}
            
        \column{0.48\textwidth}
            \textbf{Paging:}
            \begin{itemize}
                \item Membagi memori fisik menjadi frame berukuran tetap
                \item Membagi memori logis menjadi halaman dengan ukuran sama
                \item Menghilangkan external fragmentation
                \item Tetapi mungkin menyebabkan internal fragmentation
            \end{itemize}
            
            \textbf{Buddy System:}
            \begin{itemize}
                \item Algoritma alokasi yang meminimalkan fragmentasi
                \item Memecah dan menggabungkan blok memori dalam ukuran pangkat 2
                \item Memudahkan penggabungan blok memori yang berdekatan
            \end{itemize}
    \end{columns}
\end{frame}

% --- Frame: Kompaksi Memori ---
\begin{frame}{Kompaksi Memori}
    \framesubtitle{Mengatasi External Fragmentation}
    
    \begin{block}{Definisi}
        Kompaksi adalah proses menggeser isi memori untuk menggabungkan semua \textit{free space} menjadi satu blok besar.
    \end{block}
    
    \begin{columns}[T]
        \column{0.48\textwidth}
            \textbf{Sebelum Kompaksi:}
            \begin{center}
                \begin{tabular}{|c|c|c|c|c|}
                \hline
                Proses A & \cellcolor{gray!30}Hole & Proses B & \cellcolor{gray!30}Hole & Proses C \\
                \hline
                \end{tabular}
            \end{center}
            
            \textbf{Setelah Kompaksi:}
            \begin{center}
                \begin{tabular}{|c|c|c|c|}
                \hline
                Proses A & Proses B & Proses C & \cellcolor{gray!30}Hole Besar \\
                \hline
                \end{tabular}
            \end{center}
            
        \column{0.48\textwidth}
            \textbf{Tantangan Kompaksi:}
            \begin{itemize}
                \item Biaya waktu tinggi untuk memindahkan proses
                \item Membutuhkan CPU untuk operasi move
                \item Hanya mungkin jika relokasi dinamis tersedia
                \item Semua alamat dalam memori harus diperbarui
                \item Tidak praktis dilakukan terlalu sering
            \end{itemize}
    \end{columns}
\end{frame}

% --- Frame: Kerugian dari Fragmentasi ---
\begin{frame}{Kerugian dari Fragmentasi}
    \framesubtitle{Dampak pada Performa Sistem}
    
    \begin{block}{Dampak Jangka Pendek}
        \begin{itemize}
            \item Pemborosan memori fisik
            \item Permintaan alokasi gagal meskipun memori total cukup
            \item Waktu tunggu proses meningkat
            \item Peningkatan overhead sistem karena manajemen fragmen kecil
        \end{itemize}
    \end{block}
    
    \begin{alertblock}{Dampak Jangka Panjang}
        \begin{itemize}
            \item Penurunan performa sistem secara keseluruhan
            \item Peningkatan frekuensi kompaksi memori (overhead tambahan)
            \item Peningkatan aktivitas swapping ke disk
            \item Waktu tanggap sistem memburuk
            \item Throughput sistem menurun
        \end{itemize}
    \end{alertblock}
\end{frame}

% --- Frame: Studi Kasus Fragmentasi ---
\begin{frame}{Studi Kasus Fragmentasi}
    \framesubtitle{Sistem Server Web}
    
    \begin{columns}[T]
        \column{0.48\textwidth}
            \textbf{Skenario:}
            \begin{itemize}
                \item Server web dengan memori 8GB
                \item Menjalankan puluhan worker process
                \item Setiap worker process membutuhkan memori berbeda
                \item Proses dimulai dan diakhiri sesuai permintaan
                \item Sistem berjalan 24/7 tanpa restart
            \end{itemize}
            
        \column{0.48\textwidth}
            \textbf{Masalah yang Muncul:}
            \begin{itemize}
                \item Setelah beberapa hari, server mulai menolak koneksi baru
                \item Monitor memori menunjukkan 3GB memori bebas
                \item Tetapi tidak ada blok cukup besar untuk worker baru
                \item Fragmentasi eksternal menjadi bottleneck
                \item Perlu restart untuk 'reset' alokasi memori
            \end{itemize}
    \end{columns}
    
    \begin{alertblock}{Solusi}
        Implementasi sistem paging atau algoritma alokasi memori yang lebih efisien seperti buddy system akan mengurangi dampak fragmentasi.
    \end{alertblock}
\end{frame}

% --- Frame: Kesimpulan ---
\begin{frame}{Kesimpulan Fragmentasi}
    \framesubtitle{Poin-poin Kunci}
    
    \begin{itemize}
        \item Fragmentasi adalah konsekuensi tak terhindarkan dari alokasi dinamis
        \item External fragmentation: memori terbuang di antara blok-blok terpakai
        \item Internal fragmentation: memori terbuang di dalam blok terpakai
        \item 50 Percent Rule menunjukkan besarnya potensi kerugian karena fragmentasi
        \item Strategi mengatasi fragmentasi:
        \begin{itemize}
            \item Kompaksi memori
            \item Paging dan segmentasi
            \item Algoritma alokasi efisien (Buddy system)
        \end{itemize}
        \item Fragmentasi memiliki dampak signifikan pada performa sistem jangka panjang
    \end{itemize}
    
    \begin{alertblock}{Pertimbangan Desain}
        Sistem operasi modern menggabungkan berbagai teknik untuk meminimalkan fragmentasi dan mengoptimalkan penggunaan memori.
    \end{alertblock}
\end{frame}

\backmatter % generate the final slide
\end{document}
